% Source PDF: source/普通高考/2011/2011北京文.pdf, page 1, question 6.
% Original flowchart is vector/text artwork; preserved PNG remains under img/.
% Nodes: 开始; 输入A; P=1,S=1; S≤A?; P=P+1; S=S+1/P; 输出P; 结束.
% Directed edges: 开始→输入A→初始化→判定; 判定(是)→P更新→S更新→回主干入口→判定;
% 判定(否)→输出P→结束. The return arrow and output branch are independent.

\begin{tikzpicture}[
  >=Stealth,
  x=1cm,
  y=1cm,
  line width=0.45pt,
  line cap=round,
  line join=round,
  font=\normalsize,
  flowtext/.style={anchor=center, align=center,
    execute at begin node={\setlength{\parindent}{0pt}}},
  flowbox/.style={flowtext, draw, minimum height=0.68cm,
    inner xsep=4pt, inner ysep=2pt},
  terminal/.style={flowbox, rounded corners=2pt, minimum width=1.45cm},
  process/.style={flowbox, minimum width=3.25cm},
  decision/.style={flowbox, diamond, aspect=2.8, inner sep=0pt,
    minimum width=4.55cm, minimum height=1.08cm},
  io/.style={flowbox, trapezium, trapezium left angle=72,
    trapezium right angle=108, minimum height=0.78cm},
  branchlabel/.style={flowtext, inner sep=0pt},
]
  \node[terminal] (start) at (0,11.25) {开始};
  \node[process, minimum width=2.55cm] (input) at (0,9.80) {输入 $A$};
  \node[process, minimum width=4.10cm] (init) at (0,8.30) {$P=1,S=1$};
  \node[decision] (test) at (0,6.45) {$S\leq A$};
  \node[process, minimum width=3.45cm] (incp) at (0,4.60) {$P=P+1$};
  \node[process, minimum width=4.10cm, minimum height=1.20cm]
    (updates) at (0,2.65) {$S=S+\dfrac{1}{P}$};
  \node[io, minimum width=2.85cm] (output) at (4.25,4.60) {输出 $P$};
  \node[terminal] (finish) at (4.25,2.85) {结束};

  \coordinate (loopjoin) at (0,7.45);
  \coordinate (loop-left) at (-2.85,7.45);
  \coordinate (loop-under) at (0,1.45);
  \coordinate (loop-bottom) at (-2.85,1.45);
  \coordinate (output-turn) at (4.25,6.45);

  \draw[->] (start.south) -- (input.north);
  \draw[->] (input.south) -- (init.north);
  \draw (init.south) -- (loopjoin);
  \draw[->] (loopjoin) -- (test.north);

  % 是 branch: update P, update S, then return left into the main connector.
  \draw[->] (test.south) -- (incp.north);
  \node[branchlabel] at (0.38,5.58) {是};
  \draw[->] (incp.south) -- (updates.north);
  \draw (updates.south) -- (loop-under) -- (loop-bottom) -- (loop-left);
  \draw[->] (loop-left) -- (loopjoin);

  % 否 branch goes independently to output and finish.
  \draw (test.east) -- (output-turn);
  \draw[->] (output-turn) -- (output.north);
  \node[branchlabel] at (2.00,6.72) {否};
  \draw[->] (output.south) -- (finish.north);

  \path[use as bounding box] (-3.35,1.25) rectangle (5.15,11.75);
\end{tikzpicture}
